\section{The axioms, their scope, and their implications}

\begin{claim}{A.1}{Known}{The historical and modern definitions differ}
In XVI, manuscript pp.~1--3, Loc~(1) requires isomorphisms, two out of three, and the implication
\[
gf,fg\in\W\quad\Longrightarrow\quad f,g\in\W
\]
for $f:X\to Y$ and $g:Y\to X$. Loc~(1 bis) adds closure under retracts in the arrow category. Loc~(1 ter) says that $f\in\W$ whenever its image in $\Cat[\W^{-1}]$ is invertible. Loc~(2) concerns categories with a terminal object. The literal Loc~(3) is the absolute Theorem~A statement: $u:A\to B$ belongs to $\W$ if every $A/b$ is $\W$-aspherical. The relative Theorem~A in our conventions is printed as Loc~(3 bis), not Loc~(3) \citep[ms.~1--3]{AGGrothendieckXVI}.
\end{claim}
The distinction is explicit in the manuscript; its footnote~1 even proposes exchanging the two labels. Our use of ``fundamental localizer'' follows the definition requested in the prompt and \citet[Definition~3.3.2]{RootCisinski2006}. For this definition strong saturation, including closure under retracts, is a theorem \citep[Proposition~4.2.4]{RootCisinski2006}. We do not silently apply that theorem to the weaker historical definition. The minimum for the literal Loc~(1,2,3) is nevertheless also $\Winfty$: a separate theorem for weak fundamental localizers establishes this, as explained in \claimref{A.19}.

\begin{claim}{A.2}{Known}{The exact two excision axioms of XII}
Let $X$ be a small category, $U,V\subseteq X$ full sieves, $U\cup V=X$, and $S=U\cap V$. Every arrow of $X$ lies in $U$ or $V$, so the square
\[
\begin{tikzcd}
 S\ar[r,"i_1"]\ar[d,"i_2"']&U\ar[d,"j_2"]\\
 V\ar[r,"j_1"']&X
\end{tikzcd}                                                    \tag{A.MV}
\]
is a pushout and a pullback in $\Cat$. XII, manuscript pp.~75--76, calls W(7 bis) the assertion
\[
i_1\in\W\ \Longleftrightarrow\ j_1\in\W,
\qquad i_2\in\W\ \Longleftrightarrow\ j_2\in\W.
\]
The second equivalence follows from the first by exchanging $U$ and $V$. W(7), manuscript pp.~76--77, concerns a strict natural transformation $f:X_\bullet\to Y_\bullet$ between two such squares. With $0$ the intersection, $1,2$ the pieces, and $3$ the union, it says that any two of
\[
f_0\in\W,\qquad (f_1,f_2\in\W),\qquad f_3\in\W
\]
imply the third. These are quantifications over all small categories and all such maps of covers, including empty intersections and empty pieces \citep[ms.~75--77]{AGGrothendieckXII}.
\end{claim}
Opposites exchange sieves with cosieves. Since a modern fundamental localizer is invariant under opposites, either formulation gives exactly the corresponding cosieve formulation. This is category-theoretic duality; it does not exchange homotopy pushouts with homotopy pullbacks. XII's comparison with the dual derivator at manuscript pp.~79--80 is a different operation.

\begin{claim}{A.3}{Known}{Loc~(6), its three cube clauses, and Loc~(6$'$)}
XVI, manuscript pp.~16$'$--18$'$, uses the same map $f:X_\bullet\to Y_\bullet$ of MV squares and defines
\[
\begin{array}{rcl}
\mathrm{Loc}(6a)&:& f_0,f_1,f_2\in\W\ \Longrightarrow\ f_3\in\W,\\
\mathrm{Loc}(6b)&:& f_1,f_2,f_3\in\W\ \Longrightarrow\ f_0\in\W,\\
\mathrm{Loc}(6c)&:& f_0,f_3\in\W\ \Longrightarrow\ f_1,f_2\in\W.
\end{array}
\]
After footnote~13, Loc~(6) means $(6a)+(6b)$, expressly excluding $(6c)$. On the single square (A.MV), Loc~$(6'a)$ is $i_1\in\W\Rightarrow j_1\in\W$, and Loc~$(6'b)$ is $j_1\in\W\Rightarrow i_1\in\W$. Thus XII W(7 bis) is XVI Loc~$(6')$, whereas XII W(7) includes all three clauses $(6a,b,c)$.
\end{claim}
The quoted adjective ``inhabituelle'' at XII, manuscript p.~270, applies to the reverse implication. The change of conventions in XVI is substantive: one cannot replace every occurrence of W(7) by the later label Loc~(6) before accounting for retract closure.

\begin{claim}{A.4}{Known}{Cofibrations and the four structural formulations}
For a functor $i:A\to B$, let
\[
i_!:A\backslash\Cat\longrightarrow B\backslash\Cat,
\qquad (A\to X)\longmapsto(B\to B\amalg_A X).
\]
The class $\operatorname{Cof}_{\W}$ consists of the $i$ for which $i_!$ sends every morphism whose underlying functor is in $\W$ to another such morphism. XVI, manuscript pp.~25--27, states:
\[
\begin{array}{p{20mm}p{112mm}}
Loc~$(6A)$&Every $f:A\to B$ factors in $\Cat$ as $A\xrightarrow{i}M\xrightarrow{p}B$, with $i\in\operatorname{Cof}_{\W}$ and $p\in\W$.\\[3pt]
Loc~$(6A')$&Every Dwyer map belongs to $\operatorname{Cof}_{\W}$.\\[3pt]
Loc~$(6B)$&Every MV square (A.MV) is $\W$-homotopy cocartesian.\\[3pt]
Loc~$(6C)$&In every $\W$-homotopy cocartesian square (A.MV), now without any sieve or strict-pushout hypothesis, $j_1\in\W$ implies $i_1\in\W$.
\end{array}
\]
Here $\W$-homotopy cocartesian means that the canonical comparison from a specified Grothendieck model of the homotopy pushout to the fourth vertex belongs to $\W$. Under enough cofibrations this agrees with the cofibration-replacement definition of XII, manuscript pp.~261--264.
\end{claim}
A Dwyer map, in the convention used here, is an inclusion $A\hookrightarrow B$ of a full sieve such that the inclusion into the cosieve generated by $A$ has a right adjoint restricting to the identity on $A$. Its opposite is the closed variant. Closure of this class under pushouts and compositions is not closure under retracts: the latter fails \citep{AGCisinski1999}. None of the arguments below assumes retract closure of the class of Dwyer maps.

\begin{claim}{A.5}{Known}{Strict fibrations and the auxiliary classes}
For $p:E\to S$, the class $\Fib_{\W}$ is defined by strict base change:
\[
p\in\Fib_{\W}\quad\Longleftrightarrow\quad
\bigl(u:X\to Y\text{ over }S,\ u\in\W\bigr)
\Rightarrow\bigl(E\times_S X\to E\times_S Y\bigr)\in\W.
\]
We write $\W_S^u$ for morphisms over $S$ whose pullbacks along every $S'\to S$ lie in $\W$. A functor $E\to S$ itself lies in $\W^u$ when all its strict base changes lie in $\W$. The class $\mathcal F_{\W}$ consists of Grothendieck fibrations $p:E\to S$ whose reindexing functors $\alpha^*:E_t\to E_s$, for every $\alpha:s\to t$ in $S$, lie in $\W$.

XVI, manuscript pp.~27--32, calls $p$ weakly $\W$-fibrant if, for every $T'\to T$ over $S$ with $T,T'$ contractible in the category-homotopy sense, $E\times_S T'\to E\times_S T$ is in $\W$. This defines $\Fib^{\mathrm{weak}}_{\W}$. Finally, $\Parf_{\W}$ means both $\W$-smooth and $\W$-proper. The manuscript has
\[
\Parf_{\W}\subseteq\Fib^{\mathrm{weak}}_{\W},\quad
\mathcal F_{\W}\subseteq\Fib^{\mathrm{weak}}_{\W},\quad
\Fib_{\W}\subseteq\Fib^{\mathrm{weak}}_{\W}.
\]
These are different classes; a Grothendieck fibration is not, by definition, a member of $\Fib_{\W}$.
\end{claim}

\begin{claim}{A.6}{Known}{The exact fibration axioms}
XVI, manuscript pp.~38--45, defines Loc~(7) to mean Loc~$(7a)$:
\[
\mathrm{Loc}(7a):\quad \mathcal F_{\W}\subseteq\Fib_{\W}.
\]
Loc~$(7b)$ says that, whenever $p:X\to S$ and $q:Y\to S$ belong to $\mathcal F_{\W}$, a functor $f:X\to Y$ over $S$ with $f\in\W$ lies in $\W_S^u$. Equivalently in this situation, every strict fibre map $f_s:X_s\to Y_s$ lies in $\W$. Loc~$(7c)$ says
\[
\mathcal F_{\W}\cap\W\subseteq\W^u.
\]
For a reference localizer $V\subseteq\W$, Loc~$(7D_V)$ is
\[
\Parf_V\cap\W\subseteq\W^u.
\]
Loc~$(7D)$ without a subscript means $V=\W$. The weakest reference version singled out by the manuscript is $D_{W_\omega}$, with its historical minimum $W_\omega$; the modern version used in this paper is $D_{\Winfty}$.

Loc~$(7F)$ means that every functor $f:X\to Y$ admits a factorization
\[
X\xrightarrow{u}E\xrightarrow{p}Y,
\qquad u\in\W,\quad p\in\Fib_{\W}.
\]
For fundamental localizers, all $X\to e$ belong to $\Fib_{\W}$, by product stability. Thus the associated category of fibrant objects has all objects fibrant. The word ``cofibrations'' in the initial formulation of $(7F)$ on manuscript p.~39 is a variance error; the subsequent formulations explicitly say fibrations.
\end{claim}

\begin{claim}{A.7}{Known}{The four other numbered conditions used in the audit}
Loc~(4), XVI manuscript p.~8, is
\[
\coprod_{i\in I}u_i\in\W\quad\Longleftrightarrow\quad
u_i\in\W\text{ for every }i\in I
\]
for every small family, including $I=\varnothing$. Loc~(5), manuscript p.~12, is $\W\subseteq\W_0$, where $\W_0$ consists of functors inducing bijections on connected components. Loc~(8), manuscript pp.~71--72, is closure of $\W$ under all small filtered colimits in the arrow category of $\Cat$; Loc~(8 bis) asks only for sequential colimits. Loc~(9), manuscript pp.~74--75, says that $\W$ is an accessible full subcategory of the arrow category, with its inclusion accessible. It does not mean preservation of every filtered colimit.
\end{claim}
The manuscript's cardinal convention in \S5 is not the convention in which ``$\aleph_0$-accessible'' means finitely accessible. Its countable-subcategory calculation should therefore not be paraphrased as finite presentability of all weak-equivalence data. Part~C uses the modern existence of a regular accessibility cardinal.

\begin{claim}{A.8}{Proved}{Brown's lemma in the form needed in XVI 3.2.1}
Let $\mathcal M$ have an initial object $0$ and let $C,\W\subseteq\Fl(\mathcal M)$. Assume that $\W$ contains identities and satisfies two out of three; $C$ contains isomorphisms and all $0\to X$, is closed under composition and pushout, and all pushouts along maps in $C$ exist; every map factors as a map in $C$ followed by a map in $\W$; and pushouts of maps in $C\cap\W$ belong to $C\cap\W$. Then $C\subseteq\operatorname{Cof}_{\W}$. Neither strong saturation nor closure of $C$ under retracts is needed.
\end{claim}
\begin{proof}
We first prove the elementary functor lemma. If $F:\mathcal M\to\mathcal N$ sends $C\cap\W$ into a class $V$ containing identities and satisfying two out of three, it sends $\W$ into $V$. For $u:X\to Y$ in $\W$, factor
\[
[u,\id_Y]:X\amalg Y\xrightarrow{j}Z\xrightarrow{r}Y,
\qquad j\in C, r\in\W.
\]
Let $j_X:X\to Z$ and $j_Y:Y\to Z$ be the two composites with the coproduct injections. They lie in $C$ since all objects are cofibrant, and in $\W$ because $rj_X=u$ and $rj_Y=\id_Y$. Thus $F(j_X),F(j_Y)\in V$, and $F(r)F(j_Y)=\id_{F(Y)}$ implies $F(r)\in V$. Hence $F(u)\in V$.

For an object $A$, the full subcategory of $A\backslash\mathcal M$ on $A\to X$ in $C$ satisfies the same cofibrant-object axioms. Its coproduct is pushout over $A$. An absolute factorization of an under-$A$ map gives a relative factorization since composites of cofibrations are cofibrations. For any $s:A\to A'$, pushout defines a functor from this category to the corresponding category under $A'$; it preserves $C\cap\W$ by the pushout assumption. The functor lemma therefore proves that
\[
h:X\to Y\text{ under }A,\quad A\to X,A\to Y\in C,\quad h\in\W
\ \Longrightarrow\ A'\amalg_A h\in\W.                    \tag{A.8.1}
\]

Now let $i:A\to B$ be in $C$ and $u:A\to A'$ in $\W$. We must prove that $B\to B\amalg_A A'$ belongs to $\W$. The preceding factorization of $u$ writes it as a trivial cofibration followed by the retraction of a trivial cofibration. Pushout of the first factor is in $\W$ by assumption. We can therefore reduce, by an actual factorization of the pushout map and two out of three, to the case that $u$ has a section $v:A'\to A$ in $C\cap\W$, with $uv=\id_{A'}$.

Put $B_1=A\amalg_{A'}B$, where the attaching maps are $v:A'\to A$ and $iv:A'\to B$. Write $l:A\to B_1$ and $k:B\to B_1$ for the coprojections. Both $l$ and $i$ lie in $C$, and $k\in C\cap\W$. There are two functors, or morphisms in $\mathcal M$,
\[
f=[i,\id_B]:B_1\to B,\qquad
g=[ivu,\id_B]:B_1\to B.
\]
They exist because $iv$ agrees with $ivu v$, and both retract $k$, so both belong to $\W$. The map $f$ is a morphism from $l$ to $i$ under $A$. By (A.8.1), its pushout along $u$ belongs to $\W$. Associativity of ordinary pushouts identifies $A'\amalg_A B_1$ canonically with $B$. Under this identification the comparison $B_1\to A'\amalg_A B_1$ is $g$. We obtain a commutative square
\[
\begin{tikzcd}
B_1\ar[r,"f"]\ar[d,"g"']&B\ar[d,"q"]\\
B\ar[r,"u_!f"']&A'\amalg_A B,
\end{tikzcd}
\]
with $f,g,u_!f\in\W$. Consequently $q\in\W$. Stability of $C$ under pushout applies this conclusion after any initial cobase change of $i$, proving precisely the defining condition of $i\in\operatorname{Cof}_{\W}$.
\end{proof}



This is the dual of the argument in \citet[\S4, Lemmas~1--2, pp.~428--429]{AGBrown1973}; we have supplied the maps because the phrase ``Brown's lemma'' alone can conceal the relative-category step.

\begin{claim}{A.9}{Proved}{The consequences of enough cofibrations and their duals}
Let $\W$ contain identities and satisfy two out of three. Suppose pushouts exist and every morphism of $\mathcal M$ factors as a $\W$-cofibration followed by a $\W$-equivalence. Then
\[
\operatorname{Cof}_{\W}\cap\W
=\{i:\text{every pushout of }i\text{ belongs to }\W\}.
\]
Moreover, a $\W$-equivalence between objects cofibrant under a fixed $A$ stays in $\W$ after arbitrary cobase change $A\to A'$. Dually, enough $\W$-fibrations implies that a $\W$-equivalence between objects fibrant over $S$ lies in $\W_S^u$, and $\Fib_{\W}\cap\W=\W^u$.
\end{claim}
\begin{proof}
For the equality, the right side implies the left by the identity pushout and two out of three in a morphism of pushout diagrams. Conversely, let $i:A\to B$ lie in $\operatorname{Cof}_{\W}\cap\W$, and factor $s:A\to A'$ as $A\xrightarrow{j}A_0\xrightarrow{w}A'$, where $j$ is a $\W$-cofibration and $w\in\W$. The pushout $A_0\to B_0=B\amalg_A A_0$ of $i$ lies in $\W$ because $j$ preserves $\W$. The pushout $B_0\to B'=B_0\amalg_{A_0}A'$ of $w$ lies in $\W$ because $A_0\to B_0$ is a $\W$-cofibration. The commutative square with horizontal maps $A_0\to B_0$ and $A'\to B'$ then gives $A'\to B'\in\W$ by two out of three.

For the relative assertion, factor $A\to A'$ in the same way. Pushout along $j$ preserves the given map $X\to Y\in\W$. Pushout along $w$ gives equivalences on both the source and target because the respective structural maps are cofibrations. Two out of three proves the assertion. Reversing all arrows proves the two fibration statements. This is exactly the valid formal argument of XII Proposition~16, manuscript pp.~231--232, and XVI 4.4.1.
\end{proof}

\begin{claim}{A.10}{Proved}{The elementary implications between the excision clauses}
Using only two out of three and the definitions of MV squares, one has
\[
(6a)\Rightarrow(6'a),\qquad
(6b)\Rightarrow(6'b),\qquad
(6c)\Rightarrow(6'a)+(6'b).
\]
Consequently XII W(7) implies W(7 bis), as Proposition~6 at manuscript p.~78 correctly states.
\end{claim}
\begin{proof}
For the first implication map the degenerate square $(S,S,V,V)$ to $(S,U,V,X)$. Its four vertex maps are $(\id_S,i_1,\id_V,j_1)$. For the second use the map from $(S,U,V,X)$ to $(U,U,X,X)$, with vertex maps $(i_1,\id_U,j_1,\id_X)$. Both are strict maps of sieve-cover squares, and all maps have the indicated sources and targets. The first cube gives $(6c)\Rightarrow(6'b)$, and the second gives $(6c)\Rightarrow(6'a)$. Exchanging the two pieces gives the symmetric formulations. In XVI's prose after the second cube, the repeated $f_1=j_1$ must read $f_2=j_1$.
\end{proof}

\begin{claim}{A.11}{Proved}{The missing cofibration premise and its repair}
For the modern fundamental-localizer convention, Loc~$(6a)$, $(6'a)$, $(6A)$, $(6A')$, and $(6B)$ all hold. The equivalences asserted in XVI Proposition~3.6 and its footnote~19 are therefore valid in this convention. The proof gap identified in XII footnote~136 is real in the argument as originally written: termwise replacement by maps in $\W$ does not by itself show that ordinary pushout preserves them. The Dwyer pushout comparison and the gluing proof in Part~B provide that missing premise without reverse excision.
\end{claim}
\begin{proof}
Part~B constructs the specified homotopy-pushout model and proves that objectwise $\W$-equivalences of spans induce $\W$-equivalences of these models, using relative Theorem~A. For MV covers the nerve is the union of two simplicial subsets, so the comparison to the ordinary categorical union is in $\Winfty\subseteq\W$. This gives $(6a)$ and $(6B)$. For a pushout along a Dwyer map the corresponding nerve comparison is again in $\Winfty$; the same span argument proves $(6A')$. The manuscript's cylinder, formed with the product of the endpoint inclusion $\{0,2n\}\hookrightarrow\Phi_n$ with the entire source category, then gives the factorization in $(6A)$. It is not a pushout-product construction. Every comparison here is a specified functor or specified nerve comparison, not an inference from an abstract isomorphism of homotopy types.

The historically proposed route from full $(6')$ to Dwyer cofibrations is also sound once its premises are verified: apply \claimref{A.8} with $C$ the Dwyer maps, using stability under composition and pushout, the same cylinder factorization, and the preservation of trivial Dwyer maps under pushout obtained by the full MV argument. What was missing in XII was exactly the verification of these premises before ordinary pushout was used. Brown's lemma proves the implication; it does not say that an arbitrary sieve is a cofibration.
\end{proof}

\begin{claim}{A.12}{Proved}{The reversible equivalences and the role of retracts}
For a modern fundamental localizer the following are equivalent:
\[
\mathrm{Loc}(6b),\quad\mathrm{Loc}(6'b),\quad\mathrm{Loc}(6C),
\quad\mathrm{Loc}(6),\quad\mathrm{Loc}(6'),\quad\mathrm{Loc}(6c),
\quad\mathrm{W}(7).
\]
In a historical or abstract formulation where retract closure is not yet available, the proof of $(6a,b)\Rightarrow(6c)$ requires that retracts of $\W$-aspherical objects be $\W$-aspherical. This is weaker than asking that every arrow retract of every member of $\W$ belong to $\W$, but the latter suffices.
\end{claim}
\begin{proof}
The single-square reverse axiom is equivalent to detection of $\W$-equivalences by their homotopy cofibres, with the empty-source convention made explicit in Part~B. Homotopy cofibres of opposite sides of a homotopy-pushout square agree by the canonical iterated-colimit comparison. Therefore cofibres detect those two maps simultaneously, proving $(6C)$. Its restriction to MV squares gives $(6'b)$, and \claimref{A.10} gives the converse from $(6b)$.

For a map between homotopy-pushout squares, take the pointed homotopy cofibre at each of its four vertices. The resulting square is a homotopy-pushout square: form its homotopy colimit as an iterated colimit of the cube and interchange the two colimits. The constant-point square is cocartesian, so no extra term appears. Detection by cofibres turns the assertions about $f_i$ into assertions about four aspherical objects $Z_i$. If $Z_1,Z_2,Z_3$ are aspherical, reverse excision in their square implies that $Z_0$ is aspherical. This proves $(6b)$.

If instead $Z_0,Z_3$ are aspherical, their specified cone points include a section $e\to Z_0$, which is in $\W$ by two out of three with $Z_0\to e$. The map of spans $(Z_1\leftarrow e\to Z_2)\to(Z_1\leftarrow Z_0\to Z_2)$, using that section and the identity on both outside vertices, gives a specified $\W$-equivalence from the homotopy wedge to $Z_3$. The forgetful functor from pointed spaces creates this homotopy pushout because the span indexing category is contractible. Its two intermediate vertices are retracts, in the localized category, of the homotopy wedge representing the fourth vertex: collapse the other wedge summand to the chosen point. Strong saturation and closure under retracts imply that both are aspherical. Thus $(6c)$ follows. Conversely $(6c)$ implies both single-square clauses by \claimref{A.10}, and the forward clause is already available. This proves all stated equivalences. Part~F explains why modern relative Theorem~A supplies the necessary retract property automatically.
\end{proof}
The cofibres in this proof are pointed objects, even though the original categories and functors are unpointed. The point is the specified cone vertex. Forgetting it before defining the wedge retractions would omit data needed by the argument.

\begin{claim}{A.13}{Proved}{The formal fibration implication scheme}
Let $\mathcal F$ be a class of arrows in a category with the relevant pullbacks. Define
\[
\begin{array}{rcl}
\mathrm A(\mathcal F)&:&\mathcal F\subseteq\Fib_{\W},\\
\mathrm B(\mathcal F)&:&\text{every }\W\text{-map over }S\text{ between }\mathcal F\text{-objects lies in }\W_S^u,\\
\mathrm C(\mathcal F)&:&\mathcal F\cap\W\subseteq\W^u.
\end{array}
\]
If every map factors as a $\W$-map followed by a map in $\mathcal F$, then $\mathrm A(\mathcal F)\Rightarrow\mathrm B(\mathcal F)$. If identities belong to $\mathcal F$, then $\mathrm B(\mathcal F)\Rightarrow\mathrm C(\mathcal F)$. If in addition $\mathcal F$ is closed under pullback and composition and contains all terminal maps, then
\[
\mathrm A(\mathcal F)\ \Longleftrightarrow\
\mathrm B(\mathcal F)\ \Longleftrightarrow\
\mathrm C(\mathcal F).
\]
These are the precise premises of XVI Lemmas~4.4.1 and~4.4.3.
\end{claim}
\begin{proof}
The factorization premise and $\mathrm A(\mathcal F)$ give enough $\W$-fibrations. Apply the dual of \claimref{A.9} to prove $\mathrm B(\mathcal F)$. To obtain $\mathrm C(\mathcal F)$, take the target over $S$ to be $\id_S:S\to S$. Finally, under the additional hypotheses, $\mathrm C(\mathcal F)$ says that the trivial members of $\mathcal F$ are stable under pullback. The dual of \claimref{A.8} gives $\mathcal F\subseteq\Fib_{\W}$.
\end{proof}
For $V\subseteq\W$, the factorization by a $V$-equivalence followed by a $V$-perfect map is the path-category input of the manuscript. Applying \claimref{A.13} to $\Parf_V$ gives the three equivalent versions of $D_V$. In particular $D_V$ implies $(7F)$. Passing to a smaller reference class weakens each of $\mathrm A_V,\mathrm B_V,\mathrm C_V,D_V$, simply by inclusion of the classes on which the conditions are tested.

\begin{claim}{A.14}{Proved}{Replacement transfers the weak-fibration formulations}
Suppose two classes $\mathcal F,\mathcal G$ have the following comparison property: each $p:X\to S$ in $\mathcal F$ factors as
\[
X\xrightarrow{j}\overline X\xrightarrow{\overline p}S,
\qquad j\in\W_S^u,\quad\overline p\in\mathcal G.
\]
Then $\mathrm A(\mathcal G)\Rightarrow\mathrm A(\mathcal F)$ and $\mathrm C(\mathcal G)\Rightarrow\mathrm C(\mathcal F)$. If the replacements are functorial over $S$ on morphisms between objects in $\mathcal F$, then also $\mathrm B(\mathcal G)\Rightarrow\mathrm B(\mathcal F)$.
\end{claim}
\begin{proof}
Given $T'\to T$ in $\W$ over $S$, compare its pullbacks along $X\to S$ and $\overline X\to S$. The two vertical comparison maps are in $\W$ because $j$ is universally in $\W$; the bottom horizontal map is in $\W$ by $\mathrm A(\mathcal G)$. Two out of three gives the top map. For $\mathrm C$, if $p\in\W$, then $\overline p\in\W$ since $j\in\W$, so both factors remain in $\W$ after any base change. For $\mathrm B$, functoriality supplies a commutative square $j_Yf=\overline f j_X$ over $S$. If $f\in\W$, then $\overline f\in\W$, hence $\overline f\in\W_S^u$; the same square after arbitrary pullback proves $f\in\W_S^u$.
\end{proof}
This proves XVI Lemma~4.4.9 and identifies exactly why its last assertion requires more than a separate factorization for each object. In the manuscript the replacement is the functorial Grothendieck-fibration replacement $X\to\Phi_S(X)$; its universal $V$-equivalence for weakly $V$-fibrant $X\to S$ gives Corollary~4.4.10. This supplies the equivalence between using $\mathcal F_V$ and using $\Fib^{\mathrm{weak}}_V$ in each of $\mathrm A_V,\mathrm B_V,\mathrm C_V$.

\begin{claim}{A.15}{Known}{The complete list of additional arrows asserted in XVI \S4}
Besides the formal implications just proved, XVI, manuscript pp.~45--70, asserts the following. For $V\subseteq\W$ one has
\[
\mathrm A_V\Rightarrow\mathrm B_V\Rightarrow\mathrm C_V\Rightarrow D_V\Rightarrow(7F),
\]
and the same chain with $V=\W$. It asserts
\[
\begin{aligned}
\mathrm A_\W&\Longleftrightarrow
\bigl(\mathrm C_\W+\mathcal F_\W\text{ closed under composition}\bigr)
\\
&\Longleftrightarrow
\bigl(\mathrm C_\W+\Fib^{\mathrm{weak}}_\W\text{ closed under composition}\bigr).
\end{aligned}
\]
It asserts $\mathrm A_V\Leftrightarrow\mathrm B_V\Leftrightarrow\mathrm C_V$ if $\mathcal F_V$ is closed under composition, in particular if $V$ satisfies Loc~(7). Its summary 4.4.13 additionally asserts that sequential-colimit closure of $V$ makes all four reference conditions equivalent; sequential-colimit closure of both $V$ and $\W$ makes all eight conditions equivalent; and strong saturation plus Loc~(3 bis) for $\W$ gives $\mathrm B_\W\Rightarrow\mathrm A_\W$. The printed summary also says that $(7F)$ for $V$ suffices to identify the first three reference conditions; its earlier argument at 4.4.6 had used the stronger condition Loc~(7).
\end{claim}
The composition equivalences follow from \claimref{A.13}, \claimref{A.14}, and the identity $\mathcal F_\W=\mathrm{Catfib}\cap\Fib^{\mathrm{weak}}_\W$. Once $\Fib^{\mathrm{weak}}_\W\subseteq\Fib_\W$, equality holds, so composition closure follows. The reverse implications use Brown's lemma with the stated composition hypothesis. The sequential-colimit arguments use the path replacement, not reverse cofibration excision. The strong-saturation implication is attributed in the manuscript to XII Proposition~1, manuscript p.~54, using locally constant diagrams; it is not a consequence of two out of three alone. Parts~C--E settle these implications for the modern localizers considered here and distinguish $(7F)$ from them. In particular, the proposed converse $(7F)\Rightarrow D_{\Winfty}$ is false even for an accessible, strongly saturated fundamental localizer satisfying Loc~$(4,5,6a,8,9)$: $\W_0$ is a counterexample, as already recognized by XVI 4.16 (the heading following 4.4.15).

\begin{claim}{A.16}{Open}{The weaker historical axiom system is not identified with the modern one here}
This audit does not prove that the literal historical Loc~$(1,2,3)$ with absolute Theorem~A implies relative Theorem~A, Loc~$(3\mathrm{bis})$. Consequently it does not use the modern strong-saturation theorem to settle the original eight-condition implication diagram for every class satisfying only that weaker definition. The unresolved statement here is precisely that implication between axiom systems; a first nonterminal test base is $[1]$, since the relative axiom over $e$ is tautological. The minimum $W_\omega=\Winfty$ is already settled by the distinct weak-localizer theorem in \claimref{A.19}; it is not part of this residual obligation.
\end{claim}
The attempted transfer would require that a morphism of Grothendieck constructions, objectwise in the proposed weak localizer, belong to it. Absolute Theorem~A tests aspherical comma categories, not arbitrary $\W$-equivalences of fibres, and therefore does not supply this step as stated. All results in this paper that use relative Theorem~A say so through the conventions. In particular, Lead~4 itself includes Loc~(3 bis), so this unresolved historical comparison does not affect the determination requested there.

\begin{claim}{A.17}{Proved}{All eight fibration variants collapse in the modern convention}
Let $V\subseteq\W$ be modern fundamental localizers. Each of the eight conditions $\mathrm A_\W,\allowbreak\mathrm B_\W,\allowbreak\mathrm C_\W,\allowbreak D_\W,\allowbreak\mathrm A_V,\allowbreak\mathrm B_V,\allowbreak\mathrm C_V,\allowbreak D_V$ is equivalent to
\[
\W=\Winfty\qquad\text{or}\qquad\W=\Fl(\Cat).
\]
No accessibility or filtered-colimit hypothesis is needed for this conclusion. Thus every conditional arrow in \claimref{A.15} is valid for the modern convention, including the summary's $(7F)$-for-$V$ formulation, although that formulation is not what the earlier manuscript proof establishes.
\end{claim}
\begin{proof}
The formal chain in \claimref{A.13}, the replacement in \claimref{A.14}, and $\Winfty\subseteq V\subseteq\W$ imply that each condition entails $D_{\Winfty}$. \claimref{E.4} proves that this condition makes the class of $\W$-equivalences stable under homotopy pullback. If $\W$ contains a map which is not a weak homotopy equivalence, its homotopy fibres, and then successive based loop spaces of a noncontractible fibre, show that $\W$ contains a nonempty disconnected aspherical space or the map $\varnothing\to e$. Pulling back one point inclusion along a point in a different component gives $\varnothing\to e\in\W$ in the first case as well. Product stability and two out of three then make every map belong to $\W$. Otherwise minimality gives $\W=\Winfty$. For $\Winfty$, the fibration and path-replacement theorems identify weak fibrations with strict homotopical base-change maps; this gives Loc~(7), hence all eight conditions. The total localizer satisfies them directly. The proof details and all comparisons with models in $\Cat$ are in \claimref{E.5}.
\end{proof}

\begin{claim}{A.18}{Proved}{The coproduct consequences accompanying the implication diagram}
Write Loc~$(4a_0)$ for preservation of finite coproducts of equivalences, Loc~$(4a)$ for preservation of all small coproducts, and Loc~$(4b)$ for reflection of equivalences by coproducts of families with a common indexing set. The assertions surrounding XVI \S3 and used in \S4 are
\[
(6a)\Rightarrow(4a_0),\qquad
(6a)+(8)\Rightarrow(4a),\qquad
(7)\Rightarrow(4b).
\]
Consequently $(6a)+(7)+(8)$ implies Loc~$(4)$. In the modern convention $(4a)$ already follows from relative Theorem~A; no reverse excision is needed for any of these assertions.
\end{claim}
\begin{proof}
Apply $(6a)$ to the disjoint-union square with empty intersection to obtain the binary case, and iterate; the empty coproduct is $\id_\varnothing$. Every small coproduct is the filtered colimit of its finite subcoproducts, which proves the second implication. For the third, view the functor $\coprod_iX_i\to\coprod_iY_i$ over the discrete category $I$. Both structural maps to $I$ are Grothendieck fibrations with identity transition functors, hence belong to $\mathcal F_\W$. By \claimref{A.13}, Loc~$(7)$ makes the given $\W$-map universally a $\W$-equivalence over $I$. Pulling back at $i\in I$ gives $X_i\to Y_i\in\W$. Finally, the modern forward coproduct assertion is the objectwise Grothendieck-construction theorem over the discrete indexing category. This also covers arbitrarily many empty summands.
\end{proof}

\begin{claim}{A.19}{Known}{The literal historical minimum is also the weak-homotopy localizer}
Let $W_\omega$ denote the intersection of all classes satisfying the literal XVI Loc~$(1,2,3)$, where Loc~$(3)$ is absolute Theorem~A. Then
\[
W_\omega=\Winfty.
\]
This equality does not require identifying every weak fundamental localizer with a modern fundamental localizer.
\end{claim}
\begin{proof}
In \citet[\S1.4.12, p.~56; Definition~6.1.12, pp.~240--241]{RootCisinski2006}, a weak fundamental localizer is a class satisfying identities, two out of three, the split-retract cancellation rule
\[
ri=\id_A,\quad ir\in\W\quad\Longrightarrow\quad i\in\W,
\]
the terminal-object axiom, and absolute Theorem~A. Literal Loc~$(1)$ implies that cancellation rule: apply its clause about $gf$ and $fg$ with $f=i$, $g=r$. Thus every historical localizer in this intersection is a weak fundamental localizer in Cisinski's precise sense. His Theorem~6.1.18, pp.~242--243, states that every such class contains $\Winfty$. Conversely $\Winfty$ satisfies the literal historical axioms, including the cancellation rule for both composites, so the intersection is exactly $\Winfty$.

For clarity about the scope of the external theorem, its proof first establishes the weak-localizer asphericity of horn inclusions by gluing and the asphericity of their transfinite cell attachments. It then obtains all weakly contractible spaces as aspherical objects. Proposition~6.1.16 passes from those objects to all $\Winfty$-equivalences using Brown's lemma and the proper Thomason model of $\Winfty$. It does not assume relative Theorem~A for the target weak localizer. This is therefore a valid separate resolution of the historical minimum, rather than an application of the modern minimum theorem with an unchecked hypothesis.
\end{proof}
